% Generated draft: requires mathematical and source review.
\documentclass[11pt]{article}
\usepackage[margin=25mm]{geometry}
\usepackage{fontspec}
\setmainfont{Latin Modern Roman}
\usepackage{amsmath,amssymb,mathtools,unicode-math}
\setmathfont{Latin Modern Math}
\usepackage{enumitem,needspace,xurl,hyperref}
\setlength{\parindent}{0pt}
\setlength{\parskip}{5pt}
\setlength{\emergencystretch}{3em}
\newcommand{\N}{\mathbb{N}}
\newcommand{\Z}{\mathbb{Z}}
\newcommand{\Q}{\mathbb{Q}}
\newcommand{\R}{\mathbb{R}}
\newcommand{\C}{\mathbb{C}}
\newcommand{\F}{\mathbb{F}}
\newcommand{\A}{\mathbb{A}}
\newcommand{\PP}{\mathbb{P}}
\newcommand{\E}{\mathbb{E}}
\newcommand{\eps}{\varepsilon}
\newcommand{\dd}{\,\mathrm{d}}
\begin{document}
\section*{UnsolvedMath ID: 30004584}
\textbf{Source problem code:} catalogue entry
\subsection{Definitions and mathematical statement}
The source concerns nineteenth-century mathematics and culture. It does not define ``worldviews'' or specify a historical corpus, geographic scope, language, or group of historical participants. For orientation, an \emph{actors' category} denotes a category used by the historical participants themselves, whereas an \emph{observers' category} denotes a category used by those studying them. These explanatory glosses do not establish which usage the source intends or whether the two usages are mutually exclusive.

\emph{Source question.} In nineteenth-century mathematics and culture, what did ``worldviews'' mean, and should the term be treated as an actors' category or an observers' category?

This is a historical and methodological research question. The supplied record specifies no mathematical objects, hypotheses, quantifiers, or criterion of resolution from which a precise mathematical proposition could be formulated. The question is therefore retained in its supplied form; its unresolved terminology and historical scope require clarification before any stronger formalization.
\subsection{Short English statement}
The question asks what ``worldviews'' meant in nineteenth-century mathematics and culture, and whether the term belongs to historical participants' own vocabulary or to historians' analytical vocabulary. The supplied record supports a methodological question, but does not determine a mathematical proposition.
\subsection{Sources}
\begin{itemize}
\item Local UnsolvedMath registry, supplied complete record for permanent ID \(30004584\), titled ``Worldviews as Historical Categories in Mathematics,'' in the Oberwolfach Reports Open Problems collection. Exact source problem code: \texttt{catalogue entry}.
\item Supplied source citation: \emph{History of Mathematics: A Global Cultural Approach} (2021), Oberwolfach Reports, DOI \texttt{10.4171/owr/2020/41}, pp.~2021--2092. \url{https://doi.org/10.4171/owr/2020/41}. The background describes Oberwolfach Reports as documenting mathematical workshops and their open-problem discussions. The record gives proposed year 2020 and names no individual proposer.
\item Supplied literature source, labeled ``original Oberwolfach report'': \url{https://doi.org/10.4171/owr/2020/41}.
\item Literature assessment dated exactly \texttt{2026-08-22}: the supplied curation labels the question open and reports that the meaning and analytical use of worldviews in historical mathematics is explicitly posed as a methodological question, with no generally accepted resolution identified. This is supplied dated curation derived from the OWR record, not an independently verified literature assessment; time-sensitive status and cited claims require independent verification.
\item The supplied statement status is \texttt{corrected\_verified}. Its accompanying verification note states that only clear formula delimiters, notation, grammar, and typography were repaired, preserving content and scope. This records the supplied editorial qualification and does not establish independent verification here.
\end{itemize}
\end{document}
